% 18-research.tex: что отложено, о чём спрашивается и что делать не репозиторию.
\section{Будущие исследования и разработки под правилом, которое запрещает большую их часть}
\label{sec:research}

В версии 2 этот раздел был перечнем того, что предстоит построить. Здесь он быть таким не может,
потому что 13~сентября 2026 года репозиторий принял операционный договор, который меняет сам смысл
слов \emph{будущая работа}. Изменению продукта теперь нужна ровно одна основательная причина:
этого требует названная внешняя реализация; этого требует названный оператор или держатель, не
являющийся автором; этого требует находка, возникшая за пределами репозитория; либо исполнимый
контрпример показывает, что поведение, которое репозиторий уже обещает, ложно. Опрятность,
полнота, изящество и желанная новая гарантия причиной не считаются, и договор велит не
придумывать пятую причину. Всё остальное уходит в лабораторию, чья работа состоит в том, чтобы
опровергать отличительные утверждения, а не расширять архитектуру. Поэтому пункты ниже
упорядочены по тому, кому пришлось бы действовать, и всё помечено \sFut, если не сказано иное.

\figfile{stigmergy}{Предложенная колония. Небольшие узкоспециальные агенты обходят код и
архитектуру и оставляют машиночитаемые феромоны на долговечном ландшафте свидетельств:
подозрение, отклонение на границе доверия, необычный переход состояния, криптографическое
допущение, отсутствующее ребро, противоречивое утверждение, место, раз за разом связанное с
отказом. Феромоны крепнут под независимым свидетельством и выветриваются без подтверждения;
агенты следуют по усиленным тропам. Распорядитель раздаёт работу и собирает свидетельства и не
решает ничего. Уверенность приходит лишь от внешних свидетельств и никогда от согласия между
моделями. То, что пережило воспроизведение, входит в собственную петлю Поляриса как тест и
инвариант.}{fig:stigmergy}

\subsection{Рубежи, которые проходить не репозиторию}

Договор принят 13~сентября 2026 года и сам себе назначает сроки. Его девяностодневная цель наступает
12~декабря 2026 года и называет семь вещей; сводная таблица (\code{лаборатория/ВНЕШНИЕ-СУЩЕСТВИТЕЛЬНЫЕ.мд})
говорит, какие из них выполнены. Считано 23~сентября 2026 года, на десятый день:

{\footnotesize
\setlength{\tabcolsep}{4pt}\renewcommand{\arraystretch}{1.12}
\rowcolors{2}{tabstripe}{white}
\begin{xltabular}{\textwidth}{>{\RaggedRight\arraybackslash\hspace{0pt}}X >{\RaggedRight\arraybackslash\hspace{0pt}}p{0.9in} >{\RaggedRight\arraybackslash\hspace{0pt}}p{2.6in}}
\caption{Девяностодневная цель операционного договора, пункт за пунктом, против сводной таблицы на десятый день.}\label{tab:ninety-day}\\
\toprule
\rowcolor{tabheader}\thead{Договор требует} & \thead{Выполнено} & \thead{Свидетельство или чего недостаёт} \\
\midrule
\endfirsthead
\toprule
\rowcolor{tabheader}\thead{Договор требует} & \thead{Выполнено} & \thead{Свидетельство или чего недостаёт} \\
\midrule
\endhead
\bottomrule
\endlastfoot
\code{полярис-проверка} устанавливается независимо & да & в реестре ПайПИ с 15~сентября 2026 года; учение о границе ставит его туда, где нет ни строки кода Поляриса \\
Настоящее криптографическое поведение объявлено явно & да & верификатор отказывается запускаться без объявленного режима криптографии (раздел~\ref{sec:verifiers}) \\
Названный сторонний кошелёк проводит предъявление & да & Уолт.ид Воллет ППИ в2, 15~сентября 2026 года \\
Работает внешний пакет проверок соответствия ОпенАйДи & да & размещённый пакет Фонда, 15~сентября 2026 года: семь \code{ПРОЙДЕНО}, четыре \code{НА РАССМОТРЕНИИ} \\
Верификатором воспользовалась доверяющая сторона или оператор, не являющиеся автором & нет & никто; это же единственное условие между кандидатом в выпуски и \code{1.0.0} \\
Повторные предъявления извне & нет & предъявлений вне НИ: ноль \\
Внешние отказы входят в петлю разработки & да & два исправления, вызванные находками извне (раздел~\ref{sec:verifiers}) \\
\end{xltabular}
}

Пять из семи выполнены к десятому дню, и обоим оставшимся нужен человек, который не есть автор, в
чём и состоит смысл договора: ничего из оставшегося в перечне нельзя сделать, написав ещё кода. За
пределами цели лежат строки, которыми не управляет никакой срок: второй кошелёк, а также кошелёк на
родном пути, а не на пути ОИД4ВП; подача на сертификацию, которая поставила бы четыре модуля в
состоянии \code{НА РАССМОТРЕНИИ} перед рецензентом Фонда; независимое ревью безопасности, проверка
на проникновение или криптографический аудит; пилот с согласившимися пользователями; и
институционально независимые свидетели.

Две оставшиеся строки может заполнить читатель настоящего документа, и репозиторий говорит, как. Путь
постороннего (\code{документы/ПУТЬ-ПОСТОРОННЕГО.мд}) ведёт от чистой машины, опубликованного пакета и
неизменённого внешнего кошелька к принятому предъявлению примерно за десять минут, без клонирования
репозитория, и его проходят от начала до конца, прежде чем изменить. Тот, кто не является автором,
прошедший его и сообщивший о результате в трекере задач репозитория, есть первая недостающая строка;
предъявление, которым он заканчивается, есть вторая. Неудача значит столько же, сколько успех: это
внешняя находка, и каждая из двух записанных до сих пор стала исправлением. Сводная таблица заполняет
строку только именем и датой, так что сообщение есть свидетельство, лишь когда его автор готов быть
названным. У читателя без Докера есть дверь короче: описание каждого из трёх других пакетов, в реестрах начиная с их следующего выпуска,
открывается первым запуском на опубликованных в репозитории векторах, подлинном удостоверении и
подделанном, который занимает минуту от реестра, и каждый был пройден как посторонним
23~сентября 2026 года. \sExt

Договор называет и другой конец. 12~марта 2027 года, на сто восьмидесятый день, внешние поля
считываются. Если они останутся, по сути, пустыми, а внутренние версии, инварианты и архитектура
продолжат расти, ядро надлежит заморозить, архитектура национального масштаба перестаёт
расширяться, верификатор сохраняется, лишь если у него есть внешнее применение, а проект
отправляется в архив как эталонная реализация. Предмет настоящего документа может быть отправлен в
архив по собственному правилу его автора, и документ записывает это правило, а не спорит с ним.
\sExt

\subsection{Вопросы, на которые лаборатория всё ещё должна ответить}

Каждый из них есть измерение, а не возможность, и дисциплина лаборатории состоит в том, чтобы
публиковать модель угроз, противника, стенд, результат и ограничения вместе, с положительным
контролем прежде, чем засчитать хоть один нулевой результат.

\begin{description}
  \item[Остальная часть модели угроз связываемости.] Измеренный результат покрывает только
    устойчивые поля. Размер предъявления, временные характеристики, метаданные эмитента, сужающие
    множество анонимности, предметы о состоянии, подшитые к предъявлению, и структура протокола не
    измерены, а единственный канал, который противник способен прочесть, длина протокола, заведомо
    открыт в любом развёртывании, где размер предъявления меняется от держателя к держателю.
  \item[Девять дескрипторов пути совместимости.] Пакетная выдача, много удостоверений по одному
    ключу держателя на каждое, есть смягчение, принятое в экосистеме, и она не реализована; будет
    ли она принята, есть вопрос к внешней стороне, а цена её отсутствия теперь выражена числом.
  \item[Вторая половина пути миграции.] Откат и смешанное окно во время частичной миграции не
    измерены. Состояние алгоритма, читаемое верификатором, в реестре или в списке доверия, было бы
    новой продуктовой гарантией, поэтому лаборатория записала пробел и не стала его строить.
  \item[Принуждение против осведомлённого принуждающего.] Сигнал тревоги со стороны держателя,
    отрицаемая регистрация, при которой отсутствие кода принуждения невозможно доказать, и путь
    очистки записи о принуждении, каждое из этого есть изменение механизма, и оценка не предлагает
    ни одного. Прежде всего она просит измерения, которого репозиторий сделать не может: что в
    действительности наблюдает принуждающий, стоящий рядом с держателем.
  \item[Нижний предел множества анонимности.] Сообщать размер толпы рядом с вердиктом
    \code{ограничено} или отказываться закрывать эпоху ниже предела значило бы изменить то, что
    система обещает; и то и другое записано для решения развёртывающей организацией.
\end{description}

\subsection{Технологии предъявления, которые убрали бы соотнесение, отложены}

Версия 2 перечисляла пять семейств. Одно, попарный идентификатор, построено, с той границей,
которую излагает раздел~\ref{sec:privacy}. Остальные четыре отложены, а не запланированы, и
причина, по которой каждое не построено, есть причина договора, а не техническая.

\begin{description}
  \item[Ослеплённые или производные предъявления.] Держатель доказывает владение подписью над
    признаками, не раскрывая самой подписи. Классические построения опираются на допущения о
    спариваниях, которые ломает квантовый вычислитель; постквантовые аналоги есть открытая область
    исследований, и постквантовую позицию репозитория пришлось бы сохранить.
  \item[Одноразовые токены предъявления.] Одноразовые предметы, вычеканенные на каждое
    предъявление, ценой обращения к эмитенту либо заранее забранного пакета, что разменивается на
    автономную доступность и ненаблюдаемость эмитентом ровно так же, как утверждение о состоянии.
  \item[Анонимные удостоверения с избирательным раскрытием.] Раскрытие того, что признак
    удовлетворяет условию, без раскрытия самого признака. Сравнение со зрелой системой измерило
    разрыв точно: как эмитент Полярис не имеет удостоверения с признаками, к которому можно было бы
    применить условие, а как верификатор принимает заранее вычисленное логическое значение, которое
    эмитенту пришлось предусмотреть.
  \item[Родное анонимное условие.] Первый сценарий сравнения удовлетворяется раскрытым признаком,
    а не анонимным условием; это числится в сводной таблице как известное ограничение и сдвинется
    лишь тогда, когда о него споткнётся внешняя сторона.
\end{description}

\subsection{Метида: изучать систему, но никогда не людей}

Оценка, которую дорожная карта требует до появления хоть строки кода, должна была бы опровергнуть
главную гипотезу (что модель способна выводить свойства системы, не выводя прав людей) против
эталонной истины синтетической страны, сравнить каждую модель с детерминированным правилом и
спроектировать конституционный инвариант, который запрещал бы запретные применения структурно, на
манер шести проверок Афины. По договору сама оценка ждёт, пока в том, что она произвела бы, не
возникнет нужда у внешней стороны; ничто в Метиде не запланировано, и настоящий документ сохраняет
задание потому, что содержащийся в нём запрет связывает всякий будущий познающий слой, будет ли
Метида когда-нибудь построена или нет.

\subsection{Мирмекс: отладка по следам как внешний наблюдатель}

Версия 2 предлагала под рабочим именем Мирмекс систему отладки, которая жила бы вне Поляриса и
нападала бы на него независимо, по образцу того, как муравьиная колония согласует свои действия
без центрального плана: узкоспециальные агенты оставляют машиночитаемые свидетельства на общем
ландшафте, уверенность берётся лишь из воспроизведений и никогда из согласия между моделями, а
всякий воспроизводимый отказ входит в петлю раздела~\ref{sec:senses} как тест и инвариант.
Предложение остаётся в силе, а вместе с ним и предостережение: прежний аппарат, живший внутри
дерева и утверждавший притязания о себе самом, был удалён как строительные леса, так что Мирмекс
подчинялся бы противоположной дисциплине, внешней и состязательной, либо не был бы построен вовсе.

Версия 3 может добавить, что то, ради чего задумывался Мирмекс, показано и без него.
Мутационные учения суть самые дешёвые члены колонии, построенные внутри дерева: они не
рассуждают, они удаляют, обращают и ждут красного, и нашли они больше, чем любое ревью. А
наблюдатель, нашедший то, чего не могло найти ни одно учение, сертификат, который принял каждый
внутренний инструмент, был не колонией агентов, а одним кошельком, написанным посторонними.
Предложение сохранено потому, что кошелёк постороннего проходит один путь, а кольца внутри него
по-прежнему видят лишь инструменты их автора.

\subsection{Инженерные пункты, отложенные по договору}

Хеш-основанный подписант за тем же интерфейсом хранения; драйверы книги на основе цепей, чьи
интерфейсы объявлены; родные клиенты кошелька и браузерный мост; раздача состояний как
развёртывание, а не как спецификация; произведённая карта; и путь внешней реализации родного
протокола. Каждый из них есть строка дорожной карты, которую договор перенёс из \emph{следующего}
в \emph{когда это понадобится кому-то извне}, и собственная история репозитория есть довод в
пользу такого переноса: триста девяносто семь версий ушло на постройку колец, которых никто извне
не просил, а первая же пришедшая извне сторона нашла дефект за полдня.
